\documentclass[10pt]{article}
\usepackage{comment}
\def\Type{LessonCaelan}
\def\TypeNum{Lesson 17}
\def\TypeTitle{Antiderivatives and Indefinite Integrals }
\def\Semester{Fall 2021} %Not Used 
\def\Course{MATH 1190 \& 1210}
\def\Targets{  \bf\color{Fuchsia}I2 }

\usepackage{preambleDJ}

%%%%%%%%%%%% BEGIN DOCUMENT %%%%%%%%%%%%%%%
%%%%%%%%%%%% BEGIN DOCUMENT %%%%%%%%%%%%%%%
%%%%%%%%%%%% BEGIN DOCUMENT %%%%%%%%%%%%%%%
%%%%%%%%%%%% BEGIN DOCUMENT %%%%%%%%%%%%%%%
%%%%%%%%%%%% BEGIN DOCUMENT %%%%%%%%%%%%%%%

%%%%%%%%%%%% BEGIN DOCUMENT %%%%%%%%%%%%%%%
%%%%%%%%%%%% BEGIN DOCUMENT %%%%%%%%%%%%%%%
%%%%%%%%%%%% BEGIN DOCUMENT %%%%%%%%%%%%%%%
%%%%%%%%%%%% BEGIN DOCUMENT %%%%%%%%%%%%%%%
%%%%%%%%%%%% BEGIN DOCUMENT %%%%%%%%%%%%%%%

\begin{document}
\pagestyle{otherpages}
\thispagestyle{firstpage}
\vs{.5}

\textbf{Intended Learning Outcomes}
After this lesson, you will be able to 

\begin{itemize}
    \item Find antiderivatives of power and exponential functions. (\target{I3})
    %\item Use graph (and possibly additional information) of a function to deduce information about its antiderivative.
    \item Use algebraic simplification and/or properties of an indefinite integral to find the indefinite integral of a function. (\target{I3})

\end{itemize}

\vspace{-0.6cm}

\hrulefill\\

\begin{center}
\fbox{\it Basic Antidifferentiation Rules}
\end{center}

{\bf\color{blue} Fill in the blanks!} Let $F$ and $G$ be antiderivatives of $f$ and $g$, respectively, and let $k$ be a constant.

\begin{itemize}
\item {\bf Definition:} A function $F$ is called an {\bf antiderivative} of $f$ on an interval $I$ if \underline{\textcolor{blue}{$F'(x)$}} $=$ \underline{\textcolor{blue}{$f(x)$}} for all $x$ in $I$.
\item A particular antiderivative of $kf(x)$ is \underline{\textcolor{blue}{$kF(x)$}}.
\item A particular antiderivative of $f(x)\pm g(x)$ is \underline{\textcolor{blue}{$F(x)\pm G(x)$}}.
\item The {\bf indefinite integral} of a function is the most general antiderivative of that function: 
\[\dint f(x)\,dx = \underline{\textcolor{blue}{F(x)+C}}\]
\end{itemize}
\

\hrulefill

\

\begin{center}
\fbox{\it Antiderivatives of Exponential Functions}
\end{center}

{\bf\color{blue} Warm-up Exercise 1.} Fill in derivative formulas:  
$\dfrac{d}{dx}[e^x] = \underline{\textcolor{blue}{e^x}}$, \quad
$\dfrac{d}{dx}[a^x] = \underline{\textcolor{blue}{a^x \ln a}}$

{\exercise{\color{Fuchsia}I3}} Use basic antidifferentiation rules to respond to the following prompts:
\begin{itemize}
\item[(a)] Find an antiderivative of $f(x)=e^x$. Use the definition of an antiderivative to check your answer. ({\bf Hint:} you might ask yourself ``which function has a derivative of $f(x)=e^x$?")
\\
\textcolor{blue}{Solution:
 Since $f(x)=e^x$, $\dint e^x dx = e^x + C$\\
 \textbf{Check:} $F'(x) = \frac{d}{dx}(e^x) = e^x = f(x)$
 }{}\item[(b)] Find an antiderivative of $f(x)=a^x$, where $a>0$ but $a\not=1$. Use the definition of an antiderivative to check your answer. ({\bf Hint:} you might ask yourself ``which function has a derivative that contains an $a^x$ term?" to get started.)
\\

\textcolor{blue}{Solution: We have $f(x)=a^x$, $a>0, a\neq 1$. Since, $f'(x) = a^x \ln(a)$ the function whose derivative would be $a^x$ is $\frac{a^x}{\ln a}$. Hence, $\dint a^x dx = \frac{a^x}{\ln a} + C$  \\
\textbf{Check:} $F'(x) = \frac{d}{dx}(\frac{a^x}{\ln(a)}) = \frac{a^x \ln(a)}{\ln(a)} = a^x = f(x)$
}
\item[(c)] Based on your work above, fill in the following indefinite integral formulas:\\
\[\dint e^x dx = \underline{\textcolor{blue}{e^x + C}}, \quad \dint a^x dx = \underline{\textcolor{blue}{\frac{a^x}{\ln a} + C}}\]
\end{itemize}

\newpage

\begin{center}
\fbox{\it Antiderivatives of Power Functions \& Polynomials}
\end{center}
{\exercise{\color{Fuchsia}I3}} Find the following antiderivatives, using the space below the table to check your answers, as needed. ({\bf Hint}: you might ask yourself questions like ``which function has a derivative of $2x$?" when filling in the table.)
\renewcommand{\arraystretch}{2}

\begin{center}
\begin{tabular}{ | c | m{2.5in} || c | m{2.5in} | }
\hline
\textbf{Function} & \textbf{Particular antiderivative} & \textbf{Function} & \textbf{Particular antiderivative}\\
\hline
$2x$ & \textcolor{blue}{$x^2$} & $6x$ & \textcolor{blue}{$3x^2$} \\
\hline
$3x^2$ & \textcolor{blue}{$x^3$} & $x^2$ & \textcolor{blue}{$\dfrac{x^3}{3}$} \\
\hline
$\dfrac{3}{2}x^{1/2}$ & \textcolor{blue}{$x^{3/2}$} & $6x^{1/2}$ & \textcolor{blue}{$4x^{3/2}$} \\
\hline
$-x^{-2}$ & \textcolor{blue}{$x^{-1}$} & $2x^{-2}$ & \textcolor{blue}{$-2x^{-1}$} \\
\hline
$\dfrac{1}{x}$ & \textcolor{blue}{$\ln|x|$} & $\dfrac{2}{x}$ & \textcolor{blue}{$2\ln|x|$} \\
\hline
$-\dfrac{1}{x^2}$ & \textcolor{blue}{$x^{-1}$} & $\dfrac{2}{x^2}$ & \textcolor{blue}{$-2x^{-1}$} \\
\hline
$x^n$ ($n \neq -1$) & \textcolor{blue}{$\dfrac{x^{n+1}}{n+1}$} & $x^{-1}$ & \textcolor{blue}{$\ln|x|$} \\
\hline
\end{tabular}
\end{center}

Now, find particular antiderivatives:

\begin{itemize}
\item $f(x)=6x-x^{-2}+6x^{1/2}+\frac2x-\frac1{x^2}$: 
\[\underline{\textcolor{blue}{F(x) = 3x^2 + x^{-1} + 4x^{3/2} + 2\ln|x| + x^{-1}}}\]
\item $g(x)=x + x^{-3} + 7x^{1/3} - x^{-1} + \frac{1}{\sqrt[3]{x}}$: 
\[\textcolor{blue}{\underline{G(x) = \frac{x^2}{2} - \frac{1}{2} x^{-2} + \frac{21}{4}x^{4/3} - \ln|x| + \frac{3}{2}x^{2/3}}}\]
\end{itemize}

Based on your work above, fill in the following indefinite integral formulas:\\

\,\hfill
$\dint x^n\,dx = \underline{\textcolor{blue}{\frac{x^{n+1}}{n+1} + C}}$
\hfill
$\dint \frac1x\,dx = \underline{\textcolor{blue}{\ln|x| + C}}$
\hfill\,

\

\newpage

\begin{center}
\fbox{\it Summary of Indefinite Integrals}
\end{center}

{\bf\color{blue} Fill in the blanks!} Assume $k$ is any constant. Furthermore, assume $f$ has antiderivative $F$ and $g$ has antiderivative $G$. Fill in the blanks below to complete the table containing some basic indefinite integrals.


\begin{center}
{\it Integral of a Constant}\\[0.1in]
\begin{tabular}{ r c l}
$\displaystyle\int k\, dx$ & $=$ & \underline{\textcolor{blue}{$kx+C$}}
\end{tabular}
\end{center}

\


\begin{center}
{\it Integral of Power Functions}\\[0.1in]
\begin{tabular}{ r c l r c l}
$\displaystyle\int x^n\, dx$ & ${=}$ &  $\underline{\textcolor{blue}{\frac{x^{n+1}}{n+1}+C},\ n\neq -1}$ & $\displaystyle\int\frac1x\,dx$ & ${=}$ & \underline{$\textcolor{blue}{\ln|x| + C}$}
\end{tabular}
\end{center}

\

\begin{center}
{\it Integral of Exponential Functions}\\[0.1in]
\begin{tabular}{ r c l r c l}
$\displaystyle\int e^x\, dx$ & $=$ & \underline{$\textcolor{blue}{e^x + C}$} & $\displaystyle\int a^x\,dx$ & $=$ & 
$\underline{\textcolor{blue}{\frac{a^x}{\ln a} + C}}$ 
\text{ (assuming $a>0\,\&\,a\not=1$) }
\end{tabular}
\end{center}

\

\begin{center}
{\it Linearity Properties of Indefinite Integrals}\\[0.1in]
\begin{tabular}{ r c l r c l}
$\displaystyle\int \big( f(x) \pm g(x) \big)\, dx$ & $=$ & \underline{$\textcolor{blue}{\dint f(x) dx \pm \dint g(x) dx}$} $=$ & \underline{$\textcolor{blue}{F(x) \pm G(x)+C}$} \text{ (where $F'(x) = f(x)$ and $G'(x) = g(x)$) }
\end{tabular}
\end{center}

\

\begin{center}
{\it Linearity Properties of Indefinite Integrals}\\[0.1in]
\begin{tabular}{ r c l r c l}
& $\displaystyle\int kf(x) \,dx$ & $=$ & \underline{$\textcolor{blue}{k \dint f(x) dx}$} $=$ & \underline{$\textcolor{blue}{kF(x)}$} \text{ (where $F'(x) = f(x)$) }
\end{tabular}
\end{center}
  \

{\exercise{\color{Fuchsia}I3}}  Find the following indefinite integrals. ({\bf Hint:} when you can't find an answer using the rules above, try rewriting and/or simplifying algebraically first!) 

\begin{enumerate}
\item[(a)] $\dint \sqrt2 dx = \textcolor{blue}{\sqrt2 x + C}$
\item[(b)] $\dint (\sqrt[3]{u}+4u^{-1}+u^{-2}) du = \textcolor{blue}{\frac{3}{4} u^{4/3} + 4 \ln|u| - u^{-1} + C}$
\item[(c)] $\dint \frac{y^4-1}{y^2} dy$ = \textcolor{blue}{
\begin{align*}
\dint \frac{y^4-1}{y^2} dy 
&= \dint \frac{y^4}{y^2} -\frac{1}{y^2} dy\\
&= \dint y^2 -\frac{1}{y^2} dy\\
&= \frac{y^3}{3} + y^{-1} + C
\end{align*}
}
\end{enumerate}

\newpage

{\bf\color{blue}Extra Practice 1.}\ltrm{\color{Fuchsia}I3} Find the following indefinite integrals:
\begin{enumerate}
\item[(a)] $\dint \left(2 - 6t^3 + \frac{3}{t^2} - e^t \right) dt = \textcolor{blue}{2t - \frac{3}{2} t^4 - 3 t^{-1} - e^t + C}$
\item[(b)] $\dint \frac{x^3+x^2-x+1}{2x^2} dx$ = 

\textcolor{blue}{
\begin{align*}
\dint \frac{x^3+x^2-x+1}{2x^2} dx &= \dint \frac{x^3}{2x^2} + \frac{x^2}{2x^2} - \frac{x}{2x^2} + \frac{1}{2x^2} dx\\
&= \dint \frac{x}{2} + \frac{1}{2} - \frac{1}{2x} + \frac{1}{2x^2} dx\\
&= \frac{x^2}{4} + \frac{x}{2} - \frac{1}{2} \ln|x| - \frac{1}{2}x^{-1} + C
\end{align*}
}
\item[(c)] $\dint \left(z^2 + \frac1{z^2}\right) \left( z - \frac1{z} \right) dz$ = 
   \textcolor{blue}{
\begin{align*} 
\dint \left(z^2 + \frac1{z^2}\right) \left( z - \frac1{z} \right) dz &= \dint \left(z^3   - z + \frac{1}{z} - \frac{1}{z^3}\right) dz\\
&= \frac{z^4}{4} - \frac{1}{2}z^{2} + \ln|z| + \frac{1}{2z^2} + C
\end{align*}
   } \
\end{enumerate}
\begin{comment}
\newpage

\begin{itemize}
\item $\dint \frac{f(x)}{g(x)} dx = \textcolor{blue}{X}$
\item $\dint \frac{lo(x)hi'(x) - hi(x)lo'(x)}{[lo(x)]^2} dx = \textcolor{blue}{-\frac{hi(x)}{lo(x)} + C}$
\item $\dint f(x) g(x) dx = \textcolor{blue}{X}$
\item $\dint (f'(x) g(x) + f(x) g'(x)) dx = \textcolor{blue}{f(x) g(x) + C}$
\item $\dint f(g(x)) dx = \textcolor{blue}{X}$
\item $\dint f'(g(x)) g'(x) dx = \textcolor{blue}{F(g(x)) + C}$
\end{itemize}
\end{comment}
\end{document}
